% GENERATED FILE. Edit canonical lessons, then run npm run generate:books.
% canonical-source-hash: fnv1a64:f1d01f56466c002b
% canonical-lessons: TE-C108-koru, TE-C108-nirnayincu, TE-C108-karanam, TE-C108-anduvalla, TE-C108-valla

\chapter{To wish for, To decide, Reason, Therefore, Because of}
\label{ch:te-to-wish-for-to-decide-reason-therefore-because-of}
\hlchaptermodality{108}

\begin{chapteropening}
\textbf{By the end of this chapter:} \emph{I can say to wish for, to decide, a reason, therefore and because of.}

Complete the last lesson of chapter 108: I can say to wish for, to decide, a reason, therefore and because of.
\end{chapteropening}

\section[\texorpdfstring{\te{కోరు} (kōru)}{కోరు (kōru)}]{\te{కోరు} (kōru) — to wish for, to ask for}

\label{lesson:TE-C108-koru}



\begin{quote}
Before the new one: say the Telugu for an intention, then the Telugu for a plan.
\end{quote}

\subsection*{You'll want to know: \te{కోరు}}
\textbf{\te{కోరు}} — \emph{kōru} — \textquotedblleft{}to wish for, to ask for\textquotedblright{}.

\begin{grammarlens}[title={What you've built}]
One more word for this chapter.
\end{grammarlens}

\subsection*{Your turn}
\begin{itemize}
\raggedright
  \item \emph{Say it:} \emph{kōru}
  \item \emph{Say it:} \emph{kōru}, once more
  \item \emph{Say it:} say \emph{kōru}
  \item \emph{Recall:} say the Telugu for an intention, then the Telugu for a plan, then say \emph{kōru} again
\end{itemize}

\begin{tcolorbox}[breakable,colback=teal!4,colframe=teal!35!black,title={Before you move on}]
What does \te{కోరు} mean? (\textquotedblleft{}To wish for, to ask for\textquotedblright{}.) Say it once more.
\end{tcolorbox}
\section[\texorpdfstring{\te{నిర్ణయించు} (nirṇayiñcu)}{నిర్ణయించు (nirṇayiñcu)}]{\te{నిర్ణయించు} (nirṇayiñcu) — to decide}

\label{lesson:TE-C108-nirnayincu}



\begin{quote}
Before the new one: say the Telugu for a plan, then the Telugu for to wish for.
\end{quote}

\subsection*{You'll want to know: \te{నిర్ణయించు}}
\textbf{\te{నిర్ణయించు}} — \emph{nirṇayiñcu} — \textquotedblleft{}to decide\textquotedblright{}.

\begin{grammarlens}[title={What you've built}]
One more word for this chapter.
\end{grammarlens}

\subsection*{Your turn}
\begin{itemize}
\raggedright
  \item \emph{Say it:} \emph{nirṇayiñcu}
  \item \emph{Say it:} \emph{nirṇayiñcu}, once more
  \item \emph{Say it:} say \emph{nirṇayiñcu}
  \item \emph{Recall:} say the Telugu for a plan, then the Telugu for to wish for, then say \emph{nirṇayiñcu} again
\end{itemize}

\begin{tcolorbox}[breakable,colback=teal!4,colframe=teal!35!black,title={Before you move on}]
What does \te{నిర్ణయించు} mean? (\textquotedblleft{}To decide\textquotedblright{}.) Say it once more.
\end{tcolorbox}
\section[\texorpdfstring{\te{కారణం} (kāraṇaṁ)}{కారణం (kāraṇaṁ)}]{\te{కారణం} (kāraṇaṁ) — a reason}

\label{lesson:TE-C108-karanam}



\begin{quote}
Before the new one: say the Telugu for to wish for, then the Telugu for to decide.
\end{quote}

\subsection*{You'll want to know: \te{కారణం}}
\textbf{\te{కారణం}} — \emph{kāraṇaṁ} — \textquotedblleft{}a reason\textquotedblright{}.

\begin{grammarlens}[title={What you've built}]
One more word for this chapter.
\end{grammarlens}

\subsection*{Your turn}
\begin{itemize}
\raggedright
  \item \emph{Say it:} \emph{kāraṇaṁ}
  \item \emph{Say it:} \emph{kāraṇaṁ}, once more
  \item \emph{Say it:} say \emph{kāraṇaṁ}
  \item \emph{Recall:} say the Telugu for to wish for, then the Telugu for to decide, then say \emph{kāraṇaṁ} again
\end{itemize}

\begin{tcolorbox}[breakable,colback=teal!4,colframe=teal!35!black,title={Before you move on}]
What does \te{కారణం} mean? (\textquotedblleft{}A reason\textquotedblright{}.) Say it once more.
\end{tcolorbox}
\section[\texorpdfstring{\te{అందువల్ల} (anduvalla)}{అందువల్ల (anduvalla)}]{\te{అందువల్ల} (anduvalla) — therefore}

\label{lesson:TE-C108-anduvalla}



\begin{quote}
Before the new one: say the Telugu for to decide, then the Telugu for a reason.
\end{quote}

\subsection*{You'll want to know: \te{అందువల్ల}}
\textbf{\te{అందువల్ల}} — \emph{anduvalla} — \textquotedblleft{}therefore\textquotedblright{}.

\begin{grammarlens}[title={What you've built}]
One more word for this chapter.
\end{grammarlens}

\subsection*{Your turn}
\begin{itemize}
\raggedright
  \item \emph{Say it:} \emph{anduvalla}
  \item \emph{Say it:} \emph{anduvalla}, once more
  \item \emph{Say it:} say \emph{anduvalla}
  \item \emph{Recall:} say the Telugu for to decide, then the Telugu for a reason, then say \emph{anduvalla} again
\end{itemize}

\begin{tcolorbox}[breakable,colback=teal!4,colframe=teal!35!black,title={Before you move on}]
What does \te{అందువల్ల} mean? (\textquotedblleft{}Therefore\textquotedblright{}.) Say it once more.
\end{tcolorbox}
\section[\texorpdfstring{\te{వల్ల} (valla)}{వల్ల (valla)}]{\te{వల్ల} (valla) — because of}

\label{lesson:TE-C108-valla}



\begin{quote}
Before the new one: say the Telugu for a reason, then the Telugu for therefore.
\end{quote}

\subsection*{You'll want to know: \te{వల్ల}}
\textbf{\te{వల్ల}} — \emph{valla} — \textquotedblleft{}because of\textquotedblright{}.

\begin{grammarlens}[title={What you've built}]
One more word for this chapter.
\end{grammarlens}

\subsection*{Your turn}
\begin{itemize}
\raggedright
  \item \emph{Say it:} \emph{valla}
  \item \emph{Say it:} \emph{valla}, once more
  \item \emph{Say it:} say \emph{valla}
  \item \emph{Recall:} say the Telugu for a reason, then the Telugu for therefore, then say \emph{valla} again
\end{itemize}

\begin{tcolorbox}[breakable,colback=teal!4,colframe=teal!35!black,title={Before you move on}]
What does \te{వల్ల} mean? (\textquotedblleft{}Because of\textquotedblright{}.) Say it once more.
\end{tcolorbox}
